\begin{table}[ht]
\centering
\begin{tabular}{|p{0.22\textwidth}|p{0.34\textwidth}|p{0.34\textwidth}|}
\hline
\textbf{Learning Paradigm} & \textbf{Biological Example} & \textbf{Machine-Learning Counterpart} \\
\hline
Operant conditioning & Matching law in pigeons (Herrnstein, 1961) & Maximum likelihood training (strict matching) \\
\hline
Reinforcement with bias/sensitivity & Generalized matching (Baum, 1974) & KL-regularized RLHF/DPO (sensitivity parameter $\beta$) \\
\hline
Inference temperature & Response variability scaling & Logit temperature scaling in sampling \\
\hline
Classical conditioning & Pavlov’s dogs (stimulus–response) & Supervised learning (input–output mapping) \\
\hline
Hebbian learning & Synaptic plasticity & Self-supervised representation learning \\
\hline
Observational learning & Imitation in primates & In-context learning / few-shot prompting \\
\hline
Memory constraints & Working memory span, recency effects & Context window size, recency bias in attention \\
\hline
\end{tabular}
\caption{Mapping of psychological learning paradigms to machine-learning counterparts, extended to include memory constraints.}
\label{tab:paradigm_mapping}
\end{table}

